% Supplementary information: sets supplement.tex as a separate PDF that goes
% alongside every version. The main text refers into it with S: labels; it
% refers back to the main text with M: labels, read from the preprint build.
\documentclass[11pt,a4paper]{article}
\usepackage[margin=25mm]{geometry}
\usepackage[T1]{fontenc}
\usepackage[numbers,sort&compress]{natbib}
\usepackage{amsmath,amssymb}
\usepackage{graphicx}
\usepackage{booktabs}
\graphicspath{{figures/}}
% Notation shortcuts only. Every macro here must be a plain \newcommand that
% expands to standard LaTeX math, so pandoc can expand it for .docx output.
\newcommand{\R}{\mathbb{R}}
\newcommand{\E}{\mathbb{E}}
\newcommand{\dd}{\mathrm{d}}
\newcommand{\vect}[1]{\mathbf{#1}}

% Everything about the paper except the main text: title, authors, abstract
% and back-matter statements. Every version reads this file and sets it in its
% own style, so this is the only place to edit them.

\newcommand{\papertitle}{A single-source LaTeX template for computational biology papers}
\newcommand{\papershorttitle}{Single-source paper template}
\newcommand{\papershortauthors}{Author et al.}
\newcommand{\paperkeywords}{LaTeX, BibTeX, journal submission, preprints}

% Authors, in order. Mark any number of corresponding authors.
%   \paperauthor{Given}{Family}{affiliation numbers}
%   \correspondingauthor{Given}{Family}{affiliation numbers}{email}
\newcommand{\paperauthors}{%
  \correspondingauthor{First}{Author}{1}{first.author@example.org}%
  \paperauthor{Second}{Author}{1,2}%
}

%   \paperaffiliation{number}{Department}{Institution}{City}{Country}
\newcommand{\paperaffiliations}{%
  \paperaffiliation{1}{Department}{Institution}{City}{Country}%
  \paperaffiliation{2}{Department}{Institution}{City}{Country}%
}

% One paragraph, no citations: most journals forbid them in the abstract.
\newcommand{\paperabstract}{%
This document is both the example paper and the documentation for a LaTeX
template. It is written once and built as a styled preprint, a review copy
with line numbers, a supplement and journal submissions whose source is plain
LaTeX in each journal's own format. The abstract, like the title, authors and
statements, is written once in \texttt{metadata.tex}; replace this text with
your own, keeping to one paragraph without citations.}

\newcommand{\paperacknowledgements}{%
Statements such as this one are written once in \texttt{metadata.tex} and
placed wherever, and under whatever heading, each version needs them.}

\newcommand{\paperfunding}{%
Funding statement, for example: This work was supported by [funder]
[grant number].}

\newcommand{\papercompeting}{%
Competing interests statement, for example: The authors declare no competing
interests.}

\newcommand{\paperavailability}{%
Data and code availability statement, for example: Code to reproduce the
analyses is available at [repository URL].}

\newcommand{\papercontributions}{%
Author contributions statement, for example using CRediT roles.}

\bibliographystyle{unsrtnat}

\usepackage{xr}
\externaldocument[M:]{build/preprint}

\usepackage[hidelinks]{hyperref}

\renewcommand{\thesection}{S\arabic{section}}
\renewcommand{\thefigure}{S\arabic{figure}}
\renewcommand{\thetable}{S\arabic{table}}
\renewcommand{\theequation}{S\arabic{equation}}

% Authors as a plain comma-separated list; affiliations are left to the main text.
\makeatletter
\def\supp@sep{}
\newcommand{\paperauthor}[3]{\supp@sep#1 #2\def\supp@sep{, }}
\newcommand{\correspondingauthor}[4]{\paperauthor{#1}{#2}{#3}}
\makeatother

\title{Supplementary information for\\\papertitle}
\author{\paperauthors}
\date{}

\begin{document}
\maketitle

% Supplementary information: sets supplement.tex as a separate PDF that goes
% alongside every version. The main text refers into it with S: labels; it
% refers back to the main text with M: labels, read from the preprint build.
\documentclass[11pt,a4paper]{article}
\usepackage[margin=25mm]{geometry}
\usepackage[T1]{fontenc}
\usepackage[numbers,sort&compress]{natbib}
\usepackage{amsmath,amssymb}
\usepackage{graphicx}
\usepackage{booktabs}
\graphicspath{{figures/}}
% Notation shortcuts only. Every macro here must be a plain \newcommand that
% expands to standard LaTeX math, so pandoc can expand it for .docx output.
\newcommand{\R}{\mathbb{R}}
\newcommand{\E}{\mathbb{E}}
\newcommand{\dd}{\mathrm{d}}
\newcommand{\vect}[1]{\mathbf{#1}}

% Everything about the paper except the main text: title, authors, abstract
% and back-matter statements. Every version reads this file and sets it in its
% own style, so this is the only place to edit them.

\newcommand{\papertitle}{A single-source LaTeX template for computational biology papers}
\newcommand{\papershorttitle}{Single-source paper template}
\newcommand{\papershortauthors}{Author et al.}
\newcommand{\paperkeywords}{LaTeX, BibTeX, journal submission, preprints}

% Authors, in order. Mark any number of corresponding authors.
%   \paperauthor{Given}{Family}{affiliation numbers}
%   \correspondingauthor{Given}{Family}{affiliation numbers}{email}
\newcommand{\paperauthors}{%
  \correspondingauthor{First}{Author}{1}{first.author@example.org}%
  \paperauthor{Second}{Author}{1,2}%
}

%   \paperaffiliation{number}{Department}{Institution}{City}{Country}
\newcommand{\paperaffiliations}{%
  \paperaffiliation{1}{Department}{Institution}{City}{Country}%
  \paperaffiliation{2}{Department}{Institution}{City}{Country}%
}

% One paragraph, no citations: most journals forbid them in the abstract.
\newcommand{\paperabstract}{%
This document is both the example paper and the documentation for a LaTeX
template. It is written once and built as a styled preprint, a review copy
with line numbers, a supplement and journal submissions whose source is plain
LaTeX in each journal's own format. The abstract, like the title, authors and
statements, is written once in \texttt{metadata.tex}; replace this text with
your own, keeping to one paragraph without citations.}

\newcommand{\paperacknowledgements}{%
Statements such as this one are written once in \texttt{metadata.tex} and
placed wherever, and under whatever heading, each version needs them.}

\newcommand{\paperfunding}{%
Funding statement, for example: This work was supported by [funder]
[grant number].}

\newcommand{\papercompeting}{%
Competing interests statement, for example: The authors declare no competing
interests.}

\newcommand{\paperavailability}{%
Data and code availability statement, for example: Code to reproduce the
analyses is available at [repository URL].}

\newcommand{\papercontributions}{%
Author contributions statement, for example using CRediT roles.}

\bibliographystyle{unsrtnat}

\usepackage{xr}
\externaldocument[M:]{build/preprint}

\usepackage[hidelinks]{hyperref}

\renewcommand{\thesection}{S\arabic{section}}
\renewcommand{\thefigure}{S\arabic{figure}}
\renewcommand{\thetable}{S\arabic{table}}
\renewcommand{\theequation}{S\arabic{equation}}

% Authors as a plain comma-separated list; affiliations are left to the main text.
\makeatletter
\def\supp@sep{}
\newcommand{\paperauthor}[3]{\supp@sep#1 #2\def\supp@sep{, }}
\newcommand{\correspondingauthor}[4]{\paperauthor{#1}{#2}{#3}}
\makeatother

\title{Supplementary information for\\\papertitle}
\author{\paperauthors}
\date{}

\begin{document}
\maketitle

% Supplementary information: sets supplement.tex as a separate PDF that goes
% alongside every version. The main text refers into it with S: labels; it
% refers back to the main text with M: labels, read from the preprint build.
\documentclass[11pt,a4paper]{article}
\usepackage[margin=25mm]{geometry}
\usepackage[T1]{fontenc}
\usepackage[numbers,sort&compress]{natbib}
\usepackage{amsmath,amssymb}
\usepackage{graphicx}
\usepackage{booktabs}
\graphicspath{{figures/}}
% Notation shortcuts only. Every macro here must be a plain \newcommand that
% expands to standard LaTeX math, so pandoc can expand it for .docx output.
\newcommand{\R}{\mathbb{R}}
\newcommand{\E}{\mathbb{E}}
\newcommand{\dd}{\mathrm{d}}
\newcommand{\vect}[1]{\mathbf{#1}}

% Everything about the paper except the main text: title, authors, abstract
% and back-matter statements. Every version reads this file and sets it in its
% own style, so this is the only place to edit them.

\newcommand{\papertitle}{A single-source LaTeX template for computational biology papers}
\newcommand{\papershorttitle}{Single-source paper template}
\newcommand{\papershortauthors}{Author et al.}
\newcommand{\paperkeywords}{LaTeX, BibTeX, journal submission, preprints}

% Authors, in order. Mark any number of corresponding authors.
%   \paperauthor{Given}{Family}{affiliation numbers}
%   \correspondingauthor{Given}{Family}{affiliation numbers}{email}
\newcommand{\paperauthors}{%
  \correspondingauthor{First}{Author}{1}{first.author@example.org}%
  \paperauthor{Second}{Author}{1,2}%
}

%   \paperaffiliation{number}{Department}{Institution}{City}{Country}
\newcommand{\paperaffiliations}{%
  \paperaffiliation{1}{Department}{Institution}{City}{Country}%
  \paperaffiliation{2}{Department}{Institution}{City}{Country}%
}

% One paragraph, no citations: most journals forbid them in the abstract.
\newcommand{\paperabstract}{%
This document is both the example paper and the documentation for a LaTeX
template. It is written once and built as a styled preprint, a review copy
with line numbers, a supplement and journal submissions whose source is plain
LaTeX in each journal's own format. The abstract, like the title, authors and
statements, is written once in \texttt{metadata.tex}; replace this text with
your own, keeping to one paragraph without citations.}

\newcommand{\paperacknowledgements}{%
Statements such as this one are written once in \texttt{metadata.tex} and
placed wherever, and under whatever heading, each version needs them.}

\newcommand{\paperfunding}{%
Funding statement, for example: This work was supported by [funder]
[grant number].}

\newcommand{\papercompeting}{%
Competing interests statement, for example: The authors declare no competing
interests.}

\newcommand{\paperavailability}{%
Data and code availability statement, for example: Code to reproduce the
analyses is available at [repository URL].}

\newcommand{\papercontributions}{%
Author contributions statement, for example using CRediT roles.}

\bibliographystyle{unsrtnat}

\usepackage{xr}
\externaldocument[M:]{build/preprint}

\usepackage[hidelinks]{hyperref}

\renewcommand{\thesection}{S\arabic{section}}
\renewcommand{\thefigure}{S\arabic{figure}}
\renewcommand{\thetable}{S\arabic{table}}
\renewcommand{\theequation}{S\arabic{equation}}

% Authors as a plain comma-separated list; affiliations are left to the main text.
\makeatletter
\def\supp@sep{}
\newcommand{\paperauthor}[3]{\supp@sep#1 #2\def\supp@sep{, }}
\newcommand{\correspondingauthor}[4]{\paperauthor{#1}{#2}{#3}}
\makeatother

\title{Supplementary information for\\\papertitle}
\author{\paperauthors}
\date{}

\begin{document}
\maketitle

% Supplementary information: sets supplement.tex as a separate PDF that goes
% alongside every version. The main text refers into it with S: labels; it
% refers back to the main text with M: labels, read from the preprint build.
\documentclass[11pt,a4paper]{article}
\usepackage[margin=25mm]{geometry}
\usepackage[T1]{fontenc}
\usepackage[numbers,sort&compress]{natbib}
\usepackage{amsmath,amssymb}
\usepackage{graphicx}
\usepackage{booktabs}
\graphicspath{{figures/}}
\input{macros}
\input{metadata}
\bibliographystyle{unsrtnat}

\usepackage{xr}
\externaldocument[M:]{build/preprint}

\usepackage[hidelinks]{hyperref}

\renewcommand{\thesection}{S\arabic{section}}
\renewcommand{\thefigure}{S\arabic{figure}}
\renewcommand{\thetable}{S\arabic{table}}
\renewcommand{\theequation}{S\arabic{equation}}

% Authors as a plain comma-separated list; affiliations are left to the main text.
\makeatletter
\def\supp@sep{}
\newcommand{\paperauthor}[3]{\supp@sep#1 #2\def\supp@sep{, }}
\newcommand{\correspondingauthor}[4]{\paperauthor{#1}{#2}{#3}}
\makeatother

\title{Supplementary information for\\\papertitle}
\author{\paperauthors}
\date{}

\begin{document}
\maketitle

\input{supplement}

\bibliography{refs}
\end{document}


\bibliography{refs}
\end{document}


\bibliography{refs}
\end{document}


\bibliography{refs}
\end{document}
